Crawling the necessary information was the very first task for the project.
To do this, the GitHub GraphQL API was used.
The code for this milestone was split up into three Python files.
The first one was the main crawl.py file
containing the code for making the API calls to the GitHub GraphQL API endpoint.
The second one was the logger.py file to initialise the logger,
so appropriate messages could be outputted to the console.
The second and final one was the split.py file that split each individual crawled pull request into its own JSON file.
The logger.py and split.py files were relatively straight forward to code and test,
since they both performed simple functions.
However, the crawl.py was a lot more complex in nature and, thus, was more problematic to write.
This was the first problem encountered during this milestone.

In fact, there are two versions of this crawl.py file that have to be written during development.
The first version is one that does not store the ``cursor'' value of the last pull request that it fetches,
so it is not as efficient.
A ``cursor'' value is effectively the unique identifier of each pull request as it is being fetched -
it is NOT the ID of the pull request itself.
The common issue with this first version of the crawler is
that it uses up the 5,000 API calls per hour that a user gets using the GitHub GraphQL API\@.
It also does this without retrieving the necessary amount of data (in this case, over 10,000 pull requests).
As a result of this, the second version does store the ``cursor'' value of the last pull request that it fetches,
and thus is a lot more efficient with its API calls.
Only 1,000--2,000 API calls are required when using this second, improved version of the crawler,
compared to all 5,000 API calls when using the original version.

The second problem encountered during this milestone was formatting
and structuring the query string that would be sent as part of the callback to the GitHub GraphQL API endpoint.
The format of this query string has
to be very precise for the endpoint to accept it and not report back errors during compilation.
After adjusting the formatting of this query string, the crawler does eventually function as expected.

Overall, the final, improved crawler file uses the requests pip library to make a callback to the GitHub GraphQL API, so it can generate a collection of JSON files containing 10 pull requests per file.
The crawler can accommodate for a single repository as input,
as well as a dictionary of key-value pairs linking the owner of a repository with the repository's name.
As this crawl file is running,
messages are outputted to the user in the console using the logger that was initialised earlier on.
These files can then be read by the split.py file to make a JSON file for each pull request,
ready for the next milestone of the project: Database.